\subsection*{A restricted prime-discrepancy input can already carry the full zero obstruction}
\label{sec:ns43-bump-strength}
The two-dimensional test family in the channel analysis is concrete,
but its uniform discrepancy estimate must not be advertised as a
weaker arithmetic foothold without checking its strength.

\begin{proposition}[What bounded translated-bump discrepancy excludes]
\label{prop:ns43-bump-strength}
Let $\phi\in C_c^\infty(-b,b)$ be real and even, put
\[
 P(s)=\int e^{sx}\phi(x)\,dx,\qquad
 R(v)=\int\phi(x+v)\phi(x)\,dx,\qquad C=P(1/2),
\]
and define
\[
 F(u)=e^{u/2}C^2-
        \sum_{n\ge2}\frac{\Lambda(n)}{\sqrt n}R(\log n-u).
\]
If $F$ is bounded for $u\ge u_0$, every nontrivial zero $\rho$ of
$\zeta$ with $\Re\rho>1/2$ must obey $P(\rho-1/2)=0$.
There are explicit nonnegative even smooth compactly supported probes
$\phi$ for which all zeros of $P$ lie on the imaginary axis. For
such a probe, the asserted boundedness would exclude every zero off
the critical line. This is a conditional implication, not a proved
bound for $F$ and not a new RH claim. For an arbitrary preassigned
pair of bumps, the matrix input forces common transform zeros at
any off-line zero; no conclusion without that qualification is asserted.
\end{proposition}
\begin{proof}
The correlation is supported in $[-2b,2b]$ and satisfies
$\int e^{sv}R(v)\,dv=P(s)^2$ by evenness. For $\Re s>1/2$,
absolute convergence of the von Mangoldt Dirichlet series gives
\begin{equation}\label{eq:ns43-bump-laplace}
 \int_{u_0}^{\infty}e^{-su}F(u)\,du
 =\frac{C^2}{s-1/2}+P(s)^2\frac{\zeta'}{\zeta}(s+1/2)+H(s),
\end{equation}
where
\begin{align*}
 H(s)={}&C^2\frac{e^{-(s-1/2)u_0}-1}{s-1/2}\\
 &+\sum_{\log n\le u_0+2b}\frac{\Lambda(n)}{\sqrt n}
       \int_{-\infty}^{u_0}e^{-su}R(\log n-u)\,du.
\end{align*}
This correction is entire: the quotient is removable and only
finitely many compactly supported integrals occur. The usual
identity $-\zeta'/\zeta(s)=\sum\Lambda(n)n^{-s}$ is used only in
its half-plane of absolute convergence. The pole at $s=1/2$ in
\eqref{eq:ns43-bump-laplace} cancels because $C=P(1/2)$.
Boundedness of $F$ makes the left side holomorphic on $\Re s>0$.
Uniqueness of meromorphic continuation therefore forces the residue
at $s=\rho-1/2$ to vanish. That residue is
$\operatorname{mult}(\rho)P(\rho-1/2)^2$, proving the first assertion.
For a matrix of two bumps apply this argument to each diagonal;
all off-line zeros must then be common zeros of their transforms.

Here is a zero-independent choice of probe. Fix $0<B<b$ and set
$b_k=B/(\zeta(2)k^2)$. Let $\mu$ be the distribution of
$\sum_{k\ge1}U_k$, where the $U_k$ are independent uniform variables
on $[-b_k,b_k]$. The sum converges absolutely since $\sum b_k=B$.
Its measure is nonnegative, even, supported in $[-B,B]$, and has
entire bilateral transform
\[
 P_0(s)=\prod_{k\ge1}\frac{\sinh(b_ks)}{b_ks}.
\]
The product converges uniformly on compact sets, since each factor
is $1+O(b_k^2)$ there and $\sum b_k^2<\infty$. It is nonzero away
from the zeros of its factors. Those zeros are exactly imaginary
points $i\pi m/b_k$, $m\ne0$.
For every integer $M>0$,
\[
 |P_0(i\xi)|\le
 \left(\prod_{k=1}^M b_k^{-1}\right)|\xi|^{-M},
\]
since all remaining sinc factors have modulus at most one. Thus
Fourier inversion gives a $C^\infty$ density supported in $[-B,B]$.
It is nonnegative and nonzero because it is the density of $\mu$;
normalizing it in $L^2$ does not change its transform zeros. This is
the required $\phi\in C_c^\infty(-b,b)$. If an orthonormal pair is
desired, take the second even smooth bump supported in
$\{B<|x|<b\}$ and normalize it; disjoint supports give orthogonality.
For the first probe no zero with $\Re(\rho-1/2)>0$ can be masked.
The functional-equation symmetry supplies the corresponding exclusion
on the other side of the critical line, conditional on the unproved
boundedness of $F$.
\end{proof}

For the proposed matrix in the channel analysis, $F(u)=E_{11}(u/2)$
for this first probe. Hence even this restricted uniform estimate can
already contain the full off-line-zero obstruction. Nothing above
proves that estimate or asserts the converse. The useful outcome of
this check is a limit on the advertised strength of a proposed next
lemma, not a new positivity argument.
